\documentclass[11pt,a4paper]{report}

\usepackage[utf8]{inputenc}
\usepackage[T1]{fontenc}
\usepackage[french]{babel}
\usepackage[margin=2.1cm]{geometry}
\usepackage{xurl}
\usepackage{hyperref}
\usepackage{booktabs}
\usepackage{xcolor}
\usepackage{longtable}
\usepackage{enumitem}
\usepackage{fancyhdr}
\usepackage{titlesec}
\usepackage{array}
\usepackage{colortbl}
\usepackage{ragged2e}
\usepackage{listings}
\usepackage{tikz}
\usepackage{float}

% ── Palette Copifi ──
\definecolor{copifiobsidian}{HTML}{050505}
\definecolor{copifiemerald}{HTML}{10B981}
\definecolor{copifisky}{HTML}{0EA5E9}
\definecolor{copifilime}{HTML}{CCFF00}
\definecolor{copifizinc500}{HTML}{71717A}
\definecolor{copificodebg}{HTML}{F4F4F5}
\definecolor{copificodeframe}{HTML}{D4D4D8}
\definecolor{copifiwarn}{HTML}{F59E0B}
\definecolor{copifidanger}{HTML}{EF4444}

\hypersetup{
    colorlinks=true,
    linkcolor=copifiemerald,
    urlcolor=copifisky,
    pdftitle={Compte rendu technique — Copifi},
    pdfauthor={Équipe Copifi},
    pdfsubject={Inventaire complet du SaaS Copifi}
}

\pagestyle{fancy}
\fancyhf{}
\fancyhead[L]{\small\textit{Compte rendu — Copifi}}
\fancyhead[R]{\small\thepage}
\renewcommand{\headrulewidth}{0.4pt}

\titleformat{\chapter}[display]
{\normalfont\huge\bfseries\color{copifiobsidian}}{\chaptertitlename\ \thechapter}{14pt}{\Huge}

\titleformat{\section}
{\normalfont\Large\bfseries\color{copifiemerald}}
{\thesection}{1em}{}

\titleformat{\subsection}
{\normalfont\large\bfseries\color{copifiobsidian}}
{\thesubsection}{1em}{}

\lstset{
    basicstyle=\ttfamily\footnotesize,
    backgroundcolor=\color{copificodebg},
    frame=single,
    rulecolor=\color{copificodeframe},
    breaklines=true,
    columns=fullflexible,
    keepspaces=true,
}

\newcommand{\filepath}[1]{\texttt{\small #1}}
\newcommand{\feat}[1]{\textbf{#1}}
\newcommand{\statusok}{\textcolor{copifiemerald}{\textbf{Implémenté}}}
\newcommand{\statuspartial}{\textcolor{copifiwarn}{\textbf{Partiel}}}
\newcommand{\statusgap}{\textcolor{copifidanger}{\textbf{Manquant}}}

\begin{document}

% ══════════════════════════════════════════════════════════════════
\begin{titlepage}
    \centering
    \vspace*{0.8cm}
    {\Huge\bfseries\color{copifiobsidian} Compte rendu technique\par}
    \vspace{0.5cm}
    {\LARGE\color{copifiemerald} Copifi\par}
    \vspace{0.35cm}
    {\large SaaS de facturation conforme 2026 \& pilotage financier TPE/PME\par}
    \vspace{0.25cm}
    {\normalsize Inventaire détaillé du produit, de l'architecture et de l'état d'avancement\par}
    \vspace{1.2cm}
    \rule{0.82\textwidth}{1pt}
    \vspace{0.9cm}
    \begin{tabular}{>{\bfseries}r p{9.4cm}}
        Dépôt & \filepath{saas\_cfo\_digital} (GitHub) \\
        Noms alternatifs (code) & Mini CFO Digital, FinFlow (modules internes) \\
        Stack & Laravel 12, PHP 8.2+, React 18, Inertia.js, Tailwind, Vite 7 \\
        Base de données & MySQL (production Cloudways) \\
        Hébergement & Cloudways (Debian, Node 22 via nvm) \\
        Version document & 1.0 — Juillet 2026 \\
        Public visé & Direction, équipe produit, stagiaires, auditeurs techniques \\
    \end{tabular}
    \vfill
    \begin{minipage}{0.9\textwidth}
        \centering
        \small\itshape
        Ce document décrit \textbf{ce que contient le site aujourd'hui}, tel qu'implémenté dans le code source.
        Il distingue les fonctionnalités livrées, partielles et absentes.
        Les promesses marketing de la landing sont croisées avec la réalité technique.
    \end{minipage}
\end{titlepage}

\tableofcontents
\cleardoublepage

% ============================================================
\chapter{Synthèse exécutive}
% ============================================================

\section{Vision produit}

Copifi est un \textbf{SaaS B2B français} destiné aux TPE, PME, indépendants et auto-entrepreneurs.
Il combine deux piliers :

\begin{enumerate}[leftmargin=1.5em]
    \item \feat{Facturation conforme} — devis, factures, avoirs, clients, catalogue, paiements,
    génération PDF/Factur-X, transmission Plateforme Agréée (PA), e-reporting, factures fournisseurs.
    \item \feat{Pilotage financier} — saisie mensuelle de KPI (CA, charges, budget marketing, clients),
    tableaux de bord (marge, CAC, LTV), insights IA (Groq), copilote conversationnel.
\end{enumerate}

\section{Positionnement réglementaire}

Le produit anticipe la \textbf{réforme de la facturation électronique} (sept.\ 2026 réception / sept.\ 2027 émission pour les micro-entreprises) :
format Factur-X EN~16931, statuts CDAR, webhooks PA, commande planifiée d'e-reporting.

\section{État global (juillet 2026)}

\begin{center}
\begin{tabular}{@{}p{5.5cm}p{2.2cm}p{6.8cm}@{}}
\toprule
\textbf{Domaine} & \textbf{État} & \textbf{Commentaire} \\
\midrule
Landing \& marketing & \statusok & \filepath{Welcome.jsx}, chat IA public \\
Auth (email + Google) & \statusok & Laravel Breeze + Socialite \\
Abonnement Stripe & \statuspartial & 1 \texttt{STRIPE\_PRICE\_ID} ; pas d'essai 14j sans CB \\
Module facturation & \statusok & CRUD complet devis/factures/clients/catalogue \\
Factur-X + email & \statusok & PJ Factur-X à l'envoi facture \\
Intégration PA & \statuspartial & Architecture + tests ; PA réelle non validée prod \\
E-reporting & \statuspartial & Commande + service ; dépend PA \\
Pilotage KPI + IA & \statusok & Dashboard, saisie, insights Groq \\
Administration & \statusok & Suspension users, édition insights \\
Pages légales & \statusgap & Liens footer sans routes \\
Responsive mobile & \statuspartial & Desktop OK ; mobile à améliorer \\
Tests automatisés & \statuspartial & PA, compliance, auth ; peu de tests facturation \\
\bottomrule
\end{tabular}
\end{center}

\section{Chiffres clés du dépôt}

\begin{itemize}[leftmargin=1.5em]
    \item \textbf{41} migrations base de données
    \item \textbf{15} modèles Eloquent
    \item \textbf{16} services applicatifs
    \item \textbf{12} classes Support (calculs, Factur-X, compliance)
    \item \textbf{$\sim$25} contrôleurs HTTP (+ auth Breeze)
    \item \textbf{31} pages React/Inertia
    \item \textbf{$\sim$70} routes web + 1 route API PA
    \item \textbf{14} fichiers de tests PHPUnit
\end{itemize}

% ============================================================
\chapter{Architecture technique}
% ============================================================

\section{Stack logicielle}

\subsection{Backend}

\begin{center}
\begin{tabular}{@{}ll@{}}
\toprule
\textbf{Composant} & \textbf{Version (composer.json)} \\
\midrule
PHP & $\geq$ 8.2 \\
Laravel Framework & $\geq$ 12.0 \\
Inertia Laravel & $\geq$ 2.0 \\
Laravel Sanctum & $\geq$ 4.0 \\
Laravel Socialite (Google) & $\geq$ 5.27 \\
Tightenco Ziggy (routes JS) & $\geq$ 2.0 \\
atgp/factur-x & $\geq$ 3.3 \\
barryvdh/laravel-dompdf & 3.1.1 \\
\bottomrule
\end{tabular}
\end{center}

\subsection{Frontend}

\begin{center}
\begin{tabular}{@{}ll@{}}
\toprule
\textbf{Composant} & \textbf{Version (package.json)} \\
\midrule
React / React DOM & 18.2 \\
Inertia React & 2.0 \\
Vite & 7.0.7 \\
Tailwind CSS & 3.2.1 \\
Framer Motion & 12.40 \\
Recharts (graphiques) & 3.8 \\
Lucide React (icônes) & 1.16 \\
\bottomrule
\end{tabular}
\end{center}

\section{Patron architectural}

\begin{itemize}[leftmargin=1.5em]
    \item \feat{Monolithe Laravel} avec SPA React via \textbf{Inertia.js} (pas d'API REST séparée pour l'UI).
    \item \feat{Multi-tenant léger} : chaque utilisateur possède ses tiers, articles, documents (\texttt{user\_id} sur les tables facturation).
    \item \feat{Queue database} : jobs PA asynchrones (\texttt{TransmitInvoiceToPaJob}).
    \item \feat{Sessions / cache / queue} en base MySQL par défaut.
    \item \feat{Assets} compilés par Vite dans \filepath{public/build/}.
\end{itemize}

\section{Structure des répertoires principaux}

\begin{lstlisting}
app/
  Http/Controllers/     # Contrôleurs web + API PA
  Models/               # 15 modeles Eloquent
  Services/             # Logique metier
  Support/              # Calculs, Factur-X, compliance
  Jobs/                 # Transmission PA
  Mail/                 # Emails facture/devis
  Console/Commands/     # E-reporting planifie
resources/js/
  Pages/                # 31 pages Inertia
  Layouts/              # 4 layouts
  Components/           # FinFlow, Landing, Dashboard, UI
  utils/                # TVA, devises, totaux
routes/
  web.php, api.php, auth.php, console.php
database/migrations/    # 41 migrations
tests/                  # Feature + Unit
\end{lstlisting}

\section{Middleware et sécurité}

\begin{itemize}[leftmargin=1.5em]
    \item \texttt{auth}, \texttt{verified} — authentification + email vérifié
    \item \texttt{active.subscription} — abonnement Stripe \texttt{active|trialing}, compte non suspendu
    \item \texttt{admin} — rôle administrateur
    \item CSRF exempt : \filepath{/stripe/webhook}, \filepath{/api/pa/webhook}
    \item Throttle : chat landing (12/min), copilot (20/min), webhooks Stripe (60/min)
    \item Health check Laravel : \filepath{/up}
\end{itemize}

% ============================================================
\chapter{Parcours utilisateur et pages}
% ============================================================

\section{Landing et acquisition}

\begin{longtable}{@{}p{3.2cm}p{2.8cm}p{8.5cm}@{}}
\toprule
\textbf{Page} & \textbf{Route} & \textbf{Contenu} \\
\midrule
\endhead
\filepath{Welcome.jsx} & \texttt{/} &
Landing Copifi : hero, fonctionnalités, tarifs, FAQ, chat IA public (Groq), CTA inscription.
Redirige utilisateurs connectés (admin $\rightarrow$ \texttt{/admin}, abonné $\rightarrow$ \texttt{/dashboard}). \\
\filepath{TestLanding.jsx} & \texttt{/test} & Variante test landing \\
\bottomrule
\end{longtable}

\section{Authentification et compte}

Pages Breeze + Google OAuth (\filepath{Auth/}) :
Login, Register, Forgot/Reset password, Verify email, Confirm password.
Layout partagé \filepath{AuthShell.jsx} (hero masqué \texttt{< lg}).

\begin{itemize}[leftmargin=1.5em]
    \item Google : \texttt{/auth/google/redirect} et \texttt{/callback}
    \item Profil : \filepath{Profile/Edit.jsx} — mise à jour nom, email, mot de passe, suppression compte
    \item Abonnement : Stripe Checkout (\texttt{/billing/checkout}), success/cancel
\end{itemize}

\section{Pilotage financier (Mini CFO)}

\begin{longtable}{@{}p{3.5cm}p{3cm}p{7.8cm}@{}}
\toprule
\textbf{Page} & \textbf{Route} & \textbf{Contenu} \\
\midrule
\endhead
\filepath{Dashboard.jsx} & \texttt{/dashboard} &
KPIs : marge, CAC, LTV, alertes ; graphiques Recharts ; insights IA mensuels ; widget chat copilote. \\
\filepath{FinancialEntry.jsx} & \texttt{/saisie-mensuelle} &
Formulaire saisie mensuelle : CA, charges, budget marketing, nombre clients. \\
\filepath{Copilot.jsx} & \texttt{/copilote} &
Page dédiée copilote IA avec contexte financier + facturation. \\
\bottomrule
\end{longtable}

\section{Module facturation (FinFlow)}

Navigation sidebar : Vue d'ensemble, Devis, Factures, Clients, Catalogue, Paiements, Paramètres, Factures fournisseurs.

\begin{longtable}{@{}p{4cm}p{2.8cm}p{7.5cm}@{}}
\toprule
\textbf{Page} & \textbf{Route} & \textbf{Contenu} \\
\midrule
\endhead
\filepath{Facturation/Dashboard.jsx} & \texttt{/facturation} &
KPIs facturation, graphiques donut, analyse IA, timeline documents, encours. \\
\filepath{Factures/Index.jsx} & \texttt{/factures} &
Liste factures, pipeline statuts, filtres, actions PDF/envoi/paiement. \\
\filepath{FinFlow/InvoiceCreate.jsx} & \texttt{/factures/nouveau}, \texttt{/edit} &
Création/édition facture : client, lignes, remises, TVA multi-pays, devises, prestation, escompte, preview PDF. \\
\filepath{Devis/Index.jsx} & \texttt{/devis} & Liste devis + statuts \\
\filepath{Devis/Create.jsx} & \texttt{/devis/nouveau}, \texttt{/edit} &
Création devis (similaire facture) + conversion facture. \\
\filepath{FinFlow/Clients/Index.jsx} & \texttt{/clients} &
CRUD clients (tiers) : SIRET, TVA, pays, email. \\
\filepath{FinFlow/Catalogue/Index.jsx} & \texttt{/catalogue} &
CRUD articles : SKU, type, catégorie, prix HT, image. \\
\filepath{FinFlow/Payments/Index.jsx} & \texttt{/paiements} &
Suivi paiements, export CSV, retry. \\
\filepath{FinFlow/Settings/Index.jsx} & \texttt{/parametres} &
Identité entreprise, logo, SIRET, TVA, mandat facturation électronique, destinations livraison. \\
\filepath{FacturesFournisseurs/Index.jsx} & \texttt{/achats/factures-fournisseurs} &
Boîte de réception factures fournisseurs (webhook PA). \\
\bottomrule
\end{longtable}

\section{Administration}

\filepath{Admin/Dashboard.jsx} (\texttt{/admin}) : liste utilisateurs, statistiques, suspension/restauration, édition insights IA, accès dashboard utilisateur.

% ============================================================
\chapter{Routes HTTP — inventaire}
% ============================================================

\section{Routes publiques}

\texttt{GET /}, \texttt{GET /test}, \texttt{POST /landing/chat}, \texttt{GET /tracking/\{document\}/pixel}.

\section{Routes authentification (\filepath{routes/auth.php})}

Inscription, connexion, Google OAuth, reset password, vérification email, logout — standard Laravel Breeze.

\section{Routes pilotage (abonnement requis)}

\texttt{/dashboard}, \texttt{/saisie-mensuelle}, \texttt{/copilote}, \texttt{POST /dashboard/chat}, \texttt{POST /copilote/chat}, \texttt{POST /financial-records}, \texttt{GET /mes-enregistrements-financiers} (JSON).

\section{Routes facturation (abonnement requis)}

Préfixe groupé sous middleware \texttt{active.subscription} :

\begin{itemize}[leftmargin=1.5em]
    \item \texttt{/facturation} — dashboard + analyse IA
    \item \texttt{/factures/*} — CRUD, PDF, envoi, paiements, avoir
    \item \texttt{/devis/*} — CRUD, PDF, email, statuts, conversion
    \item \texttt{/clients}, \texttt{/catalogue}, \texttt{/paiements}, \texttt{/parametres}
    \item \texttt{/achats/factures-fournisseurs}
\end{itemize}

\section{Routes billing et webhooks}

\begin{itemize}[leftmargin=1.5em]
    \item Stripe Checkout : \texttt{POST/GET /billing/checkout}
    \item \texttt{POST /stripe/webhook} — événements abonnement
    \item \texttt{POST /api/pa/webhook} — statuts CDAR + factures entrantes PA
\end{itemize}

\section{Routes admin}

\texttt{/admin}, suspension/restauration users, édition insights, vue dashboard utilisateur.

% ============================================================
\chapter{Backend — modèles et domaine métier}
% ============================================================

\section{Modèles Eloquent}

\begin{longtable}{@{}p{3.2cm}p{3cm}p{7.3cm}@{}}
\toprule
\textbf{Modèle} & \textbf{Table} & \textbf{Rôle} \\
\midrule
\endhead
\texttt{User} & \texttt{users} & Compte, rôle, Stripe, Google, suspension \\
\texttt{Tier} & \texttt{tiers} & Client ou prospect (SIRET, TVA, pays) \\
\texttt{Document} & \texttt{documents} & Devis, facture ou avoir \\
\texttt{LigneDocument} & \texttt{ligne\_documents} & Lignes de document \\
\texttt{Article} & \texttt{articles} & Catalogue produits/services \\
\texttt{Payment} & \texttt{payments} & Paiements clients \\
\texttt{CompanySetting} & \texttt{company\_settings} & Identité émetteur \\
\texttt{DeliveryDestination} & \texttt{delivery\_destinations} & Frais port par destination \\
\texttt{FinancialRecord} & \texttt{financial\_records} & Saisie KPI mensuelle \\
\texttt{AiInsight} & \texttt{ai\_insights} & Insight IA mensuel \\
\texttt{DocumentEvent} & \texttt{document\_events} & Timeline audit document \\
\texttt{FactureFournisseur} & \texttt{factures\_fournisseurs} & Factures reçues PA \\
\texttt{AdminAuditLog} & \texttt{admin\_audit\_logs} & Journal actions admin \\
\texttt{ProcessedStripeEvent} & \texttt{processed\_stripe\_events} & Idempotence Stripe \\
\texttt{ProcessedPaWebhook} & \texttt{processed\_pa\_webhooks} & Idempotence PA \\
\bottomrule
\end{longtable}

\section{Document — cycle de vie}

\textbf{Types :} \texttt{devis}, \texttt{facture}, \texttt{avoir}.

\textbf{Statuts devis :} brouillon, envoyé, accepté, rejeté, expiré, converti.

\textbf{Statuts facture :} brouillon, envoyé, partiellement payé, payé, en retard, annulé.

\textbf{Statuts CDAR (PA) :} déposée, rejetée, refusée, suspendue, encaissée, approuvée.

\textbf{Relations :} document parent (devis $\rightarrow$ facture), lignes, paiements, événements, \texttt{pa\_document\_id}.

% ============================================================
\chapter{Services et logique métier}
% ============================================================

\section{Services applicatifs (\filepath{app/Services/})}

\begin{longtable}{@{}p{4.2cm}p{9.3cm}@{}}
\toprule
\textbf{Service} & \textbf{Responsabilité} \\
\midrule
\endhead
\texttt{FinancialService} & Calcul marge, CAC, LTV, alertes dashboard pilotage \\
\texttt{FinancialAnalysisService} & Analyse IA des données facturation \\
\texttt{AiInsightService} & Génération insights mensuels Groq \\
\texttt{GroqApiService} & Client HTTP API Groq LLM \\
\texttt{DashboardChatService} & Chat copilote (contexte user + KPI + facturation) \\
\texttt{LandingChatService} & Chat public landing (FAQ produit) \\
\texttt{CopilotContextService} & Construction contexte structuré copilote \\
\texttt{DocumentPdfService} & PDF DomPDF + Factur-X assemblé \\
\texttt{DocumentMailerService} & Envoi email facture \\
\texttt{DevisMailerService} & Envoi email devis \\
\texttt{ReferenceGeneratorService} & Références DEV-/FAC-/AVO- \\
\texttt{CompanySettingsService} & Paramètres entreprise frontend/PDF \\
\texttt{DeliveryDestinationService} & Destinations livraison et frais port \\
\texttt{DocumentEventRecorder} & Journal événements document \\
\texttt{RestPaApiClient} & Client REST PA (production) \\
\texttt{GenericPaApiClient} & Simulation PA locale \\
\bottomrule
\end{longtable}

\section{Classes Support (\filepath{app/Support/})}

\begin{itemize}[leftmargin=1.5em]
    \item \texttt{DocumentTotals}, \texttt{LigneAmounts} — calculs HT/TVA/TTC, remises ligne
    \item \texttt{FinancialDiscount} — escompte pour paiement anticipé
    \item \texttt{TaxRateResolver} — TVA par pays (\filepath{config/taxes.php})
    \item \texttt{ExchangeRateResolver} — devises (\filepath{config/currencies.php})
    \item \texttt{DocumentRegulatoryFields} — champs réglementaires documents
    \item \texttt{DocumentPrestation} — type prestation, frais port
    \item \texttt{ArticleCatalog} — articles disponibles pour lignes
    \item \texttt{InvoiceComplianceValidator} — contrôles pré-envoi (SIRET, catégorie opération…)
    \item \texttt{FacturXCiiGenerator} — XML CII EN~16931
    \item \texttt{FacturXAssembler} — PDF/A-3 + XML embarqué (atgp/factur-x)
    \item \texttt{EReportingService} — collecte transactions B2C/international pour PA
\end{itemize}

\section{Jobs, emails, commandes}

\begin{itemize}[leftmargin=1.5em]
    \item \textbf{Job} \texttt{TransmitInvoiceToPaJob} — soumission Factur-X à la PA (5 tentatives, backoff)
    \item \textbf{Mail} \texttt{DocumentSentMail} — facture + PJ Factur-X
    \item \textbf{Mail} \texttt{DevisSentMail} — devis + PJ PDF
    \item \textbf{Commande} \texttt{finflow:process-ereporting} — planifiée daily 02:00 (\filepath{routes/console.php})
\end{itemize}

% ============================================================
\chapter{Fonctionnalités détaillées}
% ============================================================

\section{Facturation et documents}

\begin{itemize}[leftmargin=1.5em]
    \item CRUD \textbf{devis} avec statuts, envoi email, conversion en facture
    \item CRUD \textbf{factures} avec validation compliance, PDF, envoi email Factur-X
    \item Génération \textbf{avoir} depuis facture
    \item \textbf{Paiements} rattachés aux factures (partiel/total, escompte financier)
    \item \textbf{Remises} par ligne (pourcentage ou montant)
    \item \textbf{Multi-devises} avec taux fixes configurables
    \item \textbf{Multi-TVA} selon pays client
    \item \textbf{Catégories d'opération} (biens/services) — conformité facture électronique
    \item \textbf{Timeline} événements par document (\filepath{DocumentTimeline.jsx})
    \item \textbf{Tracking ouverture} email (pixel \texttt{/tracking/\{id\}/pixel})
    \item Export CSV catalogue et paiements
\end{itemize}

\section{Conformité et PA}

\begin{itemize}[leftmargin=1.5em]
    \item Format \textbf{Factur-X} (PDF + XML CII) via package \texttt{atgp/factur-x}
    \item Envoi email avec pièce jointe Factur-X (\filepath{DocumentSentMail})
    \item Transmission PA asynchrone si \texttt{electronic\_invoicing\_active}
    \item Webhook PA : mise à jour statut CDAR + réception factures fournisseurs
    \item E-reporting : commande batch + service dédié
    \item Mandat facturation électronique (modal + date acceptation)
    \item Validateur compliance avant émission (\filepath{InvoiceComplianceValidator})
\end{itemize}

\section{Intelligence artificielle (Groq)}

\begin{itemize}[leftmargin=1.5em]
    \item Chat \textbf{landing} public (limité, sans auth)
    \item Chat \textbf{dashboard} et page \textbf{copilote} (contexte utilisateur)
    \item \textbf{Insights mensuels} générés et éditables par admin
    \item \textbf{Analyse financière} facturation (endpoint JSON \texttt{/facturation/analyze})
    \item Modèle par défaut : \texttt{llama-3.1-8b-instant}
\end{itemize}

\section{Abonnement et monétisation}

\begin{itemize}[leftmargin=1.5em]
    \item Stripe Checkout pour souscription
    \item Webhook Stripe : activation/suspension \texttt{stripe\_status}
    \item Middleware \texttt{active.subscription} sur pilotage et facturation
    \item Landing affiche 2 plans ; code n'expose qu'un \texttt{STRIPE\_PRICE\_ID}
\end{itemize}

\section{Administration}

\begin{itemize}[leftmargin=1.5em]
    \item Liste utilisateurs avec statut abonnement
    \item Suspension / restauration compte + audit log
    \item Consultation dashboard d'un utilisateur
    \item Édition manuelle des insights IA
\end{itemize}

% ============================================================
\chapter{Frontend — composants et design}
% ============================================================

\section{Layouts}

\begin{itemize}[leftmargin=1.5em]
    \item \filepath{AppDashboardLayout.jsx} — shell principal Copifi, sidebar, menu mobile drawer (\texttt{lg:})
    \item \filepath{FacturationLayout.jsx} — en-tête page + flash messages
    \item \filepath{GuestLayout.jsx} — pages invité
    \item \filepath{TestDashboardLayout.jsx} — prototype (sidebar non responsive)
\end{itemize}

\section{Composants FinFlow (facturation)}

\texttt{DashboardHome}, \texttt{DashboardKpiCards}, \texttt{KPICard}, \texttt{DonutChartCard}, \texttt{FinancialAnalysisCard}, \texttt{DocumentTimeline}, \texttt{PaymentModal}, \texttt{ModalMandatFacturation}, \texttt{FlashMessages}, logos (\texttt{FinFlowAwwwardsLogo}, \texttt{CompanyBrandLogo}), \texttt{UserAvatar}.

\section{Composants transverses}

\texttt{CfoPageShell}, \texttt{CfoPageBackground}, \texttt{LandingChatWidget}, \texttt{DashboardChatWidget}, composants UI auth (\texttt{login-hero-panel}, \texttt{login-with-listed-provider}).

\section{Identité visuelle}

\begin{itemize}[leftmargin=1.5em]
    \item Thème \textbf{dark fintech} : obsidian, emerald (\texttt{\#10B981}), sky, lime accent
    \item Typographies : Inter (corps), Syne (display)
    \item Logo : \filepath{public/images/finflow-logo.png}, composants animés React
    \item Marque UI : \textbf{Copifi} ; legacy code : FinFlow, Mini CFO Digital
\end{itemize}

\section{Utilitaires JavaScript}

\filepath{utils/taxRates.js}, \filepath{currency.js}, \filepath{documentTotals.js}, \filepath{ligneAmounts.js}, \filepath{financialDiscount.js}, \filepath{company.js}, \filepath{user.js}.

Hook : \filepath{hooks/useClientVatRate.js}.

% ============================================================
\chapter{Base de données}
% ============================================================

\section{Tables principales (41 migrations)}

\subsection{Noyau Laravel}
\texttt{users}, \texttt{password\_reset\_tokens}, \texttt{sessions}, \texttt{cache}, \texttt{cache\_locks}, \texttt{jobs}, \texttt{job\_batches}, \texttt{failed\_jobs}.

\subsection{Facturation}
\texttt{tiers}, \texttt{articles}, \texttt{documents}, \texttt{ligne\_documents}, \texttt{payments}, \texttt{company\_settings}, \texttt{delivery\_destinations}, \texttt{document\_events}, \texttt{factures\_fournisseurs}.

\subsection{Pilotage \& admin}
\texttt{financial\_records}, \texttt{ai\_insights}, \texttt{admin\_audit\_logs}.

\subsection{Intégrations}
\texttt{processed\_stripe\_events}, \texttt{processed\_pa\_webhooks}.

\section{Seeders}

\begin{itemize}[leftmargin=1.5em]
    \item \filepath{DatabaseSeeder.php} — admin, utilisateurs démo, records financiers
    \item \filepath{FacturationDemoSeeder.php} — jeu de données facturation démo
\end{itemize}

% ============================================================
\chapter{Configuration et variables d'environnement}
% ============================================================

\section{Fichiers config (\filepath{config/})}

\texttt{app.php}, \texttt{auth.php}, \texttt{database.php}, \texttt{mail.php}, \texttt{queue.php}, \texttt{session.php}, \texttt{cache.php}, \texttt{filesystems.php}, \texttt{services.php} (Stripe, Google, Groq, PA), \texttt{company.php}, \texttt{taxes.php}, \texttt{currencies.php}.

\section{Variables \filepath{.env} essentielles}

\begin{longtable}{@{}p{4.5cm}p{8.8cm}@{}}
\toprule
\textbf{Variable} & \textbf{Usage} \\
\midrule
\endhead
\texttt{APP\_*} & Nom, URL, locale (fr), debug \\
\texttt{DB\_*} & MySQL \\
\texttt{MAIL\_*} & SMTP production (défaut : log) \\
\texttt{STRIPE\_*} & Clés, price ID, webhook secret \\
\texttt{GOOGLE\_*} & OAuth connexion Google \\
\texttt{GROQ\_*} & API IA chat + insights \\
\texttt{PA\_API\_*} & URL, clé, secret webhook PA \\
\texttt{QUEUE\_CONNECTION} & database (recommandé redis en prod lourde) \\
\bottomrule
\end{longtable}

\section{Données partagées Inertia}

Via \filepath{HandleInertiaRequests.php} : \texttt{auth.user}, \texttt{tax\_rates}, \texttt{currencies}, \texttt{flash}, \texttt{company}, \texttt{delivery\_destinations}.

% ============================================================
\chapter{Templates email et PDF}
% ============================================================

\begin{itemize}[leftmargin=1.5em]
    \item \filepath{resources/views/pdf/facture.blade.php} — template PDF facture/devis/avoir
    \item \filepath{resources/views/emails/document-sent.blade.php} — email facture (mention Factur-X)
    \item \filepath{resources/views/emails/devis-sent.blade.php} — email devis
    \item \filepath{resources/views/app.blade.php} — shell HTML Inertia + Vite
\end{itemize}

% ============================================================
\chapter{Tests automatisés}
% ============================================================

\section{Couverture actuelle}

\begin{itemize}[leftmargin=1.5em]
    \item \textbf{Auth} : login, register, reset, Google OAuth, vérification email
    \item \textbf{Profile} : CRUD profil
    \item \textbf{PA} : job transmission, webhooks statut/fournisseur, e-reporting, idempotence
    \item \textbf{Email} : pièce jointe Factur-X (\filepath{DocumentSentMailTest})
    \item \textbf{Unit} : \texttt{FinancialService}, \texttt{InvoiceComplianceValidator}
\end{itemize}

\section{Non couvert (gaps tests)}

CRUD devis/factures E2E, webhooks Stripe, rendu PDF, chat Groq, pages React, exports CSV.

% ============================================================
\chapter{Déploiement et exploitation}
% ============================================================

\section{Environnement production (Cloudways)}

\begin{itemize}[leftmargin=1.5em]
    \item Application : \texttt{hkubebmsdw} sur serveur Debian
    \item Webroot : \filepath{public\_html/public}
    \item PHP 8.2+, Composer \texttt{--no-dev}, migrations, cache config/routes/views
    \item Node 22 via nvm (\texttt{source \textasciitilde/.profile}) pour \texttt{npm run build}
    \item Cron : \texttt{schedule:run} + \texttt{queue:work} chaque minute
\end{itemize}

\section{Commandes opérationnelles}

\begin{lstlisting}[language=bash]
composer install --no-dev --optimize-autoloader
npm ci && npm run build
php artisan migrate --force
php artisan storage:link
php artisan config:cache && php artisan route:cache
php artisan finflow:process-ereporting   # manuel ou cron 02:00
\end{lstlisting}

\section{Webhooks à configurer}

\begin{itemize}[leftmargin=1.5em]
    \item Stripe $\rightarrow$ \texttt{https://domaine/stripe/webhook}
    \item PA $\rightarrow$ \texttt{https://domaine/api/pa/webhook} (header \texttt{X-PA-Signature})
\end{itemize}

% ============================================================
\chapter{Lacunes, risques et feuille de route}
% ============================================================

\section{Éléments manquants ou incomplets}

\begin{longtable}{@{}p{4.5cm}p{8.8cm}@{}}
\toprule
\textbf{Sujet} & \textbf{Détail} \\
\midrule
\endhead
Pages légales &
Liens \texttt{/mentions-legales}, \texttt{/cgv}, confidentialité dans landing/auth — \textbf{aucune route}. \\
PA production &
Client REST branché ; endpoints placeholder ; non validé contre PA agréée réelle. \\
Archivage 10 ans &
Promis landing ; non implémenté en code. \\
Signature en ligne &
Champ BDD + UI ; pas de workflow signature. \\
2FA &
Paramètres : « à venir ». \\
Responsive mobile &
Desktop abouti ; mobile/tablette à traiter (brief stagiaire existant). \\
Naming &
Copifi / FinFlow / Mini CFO Digital coexistent. \\
Essai gratuit 14j &
Landing ; pas implémenté Stripe trial sans CB. \\
Plans Stripe &
2 plans landing ; 1 seul \texttt{PRICE\_ID} en env. \\
Taux de change &
Config statique, pas de API live. \\
Tests &
Couverture partielle ; pas de CI/CD documentée. \\
\bottomrule
\end{longtable}

\section{Recommandations prioritaires}

\begin{enumerate}[leftmargin=1.5em]
    \item Publier pages légales (RGPD, CGV, mentions)
    \item Brancher PA agréée réelle + recette bout-en-bout
    \item Finaliser responsive (voir \filepath{brief-stagiaire-responsive-copifi.tex})
    \item Aligner branding et variables \texttt{APP\_NAME}
    \item Compléter tests facturation + Stripe
    \item Configurer SMTP production et queue worker stable
\end{enumerate}

% ============================================================
\chapter{Annexes}
% ============================================================

\section{Documentation interne existante}

\begin{itemize}[leftmargin=1.5em]
    \item \filepath{docs/brief-identite-visuelle-copifi.tex} — brief designer
    \item \filepath{docs/brief-stagiaire-responsive-copifi.tex} — mission responsive
    \item \filepath{docs/catalogue-fonctionnalites.tex} — catalogue fonctionnalités
    \item \filepath{docs/etat-recette-finflow.tex} — état recette / gaps prod
    \item \filepath{docs/brief-creatif-logo-mini-cfo-digital.md} — brief logo legacy
\end{itemize}

\section{Documentation réglementaire (contexte produit)}

Réforme facturation électronique France :
\begin{itemize}[leftmargin=1.5em]
    \item 1\er{} septembre 2026 — réception e-factures (toutes entreprises)
    \item 1\er{} septembre 2027 — émission + e-reporting (micro-entreprises)
    \item Formats : Factur-X, UBL, CII via Plateforme Agréée
\end{itemize}

\section{Glossaire}

\begin{description}[style=nextline, leftmargin=2cm, labelwidth=2.2cm]
    \item[Copifi] Nom produit affiché dans l'UI (2026).
    \item[FinFlow] Nom interne module facturation (code, composants).
    \item[Factur-X] PDF/A-3 avec XML CII EN~16931 embarqué.
    \item[PA] Plateforme Agréée — opérateur certifié DGFiP pour e-facturation.
    \item[CDAR] Cycle de vie statut facture côté PA (déposée, approuvée…).
    \item[E-reporting] Transmission données transactions (B2C, export) à l'administration.
    \item[Inertia.js] Pont Laravel $\leftrightarrow$ React sans API REST dédiée.
    \item[Tier] Client ou prospect dans le modèle métier.
\end{description}

\vfill
\begin{center}
\small\color{copifizinc500}
Document généré à partir de l'inventaire code source — Juillet 2026\\
Compilation : \texttt{pdflatex compte-rendu-copifi.tex} (2 passes)
\end{center}

\end{document}
